\documentclass[a4paper,landscape]{article}
\usepackage[utf8]{inputenc}
\usepackage{multicol}
\usepackage{geometry}
\geometry{margin=1.5cm}
\setlength{\columnsep}{1.2cm}
\pagestyle{empty}

\begin{document}

\begin{multicols}{3}

% V = pi r^2 h
\textbf{V = \pi r$^2$ h} \\
V = Volume do cilindro (m$^3$) \\
$\pi$ = 3,1415 \\
r = raio (m) \\
h = altura (m) \\
\vspace{1em}

% Q = m c Delta T
\textbf{Q = m c \Delta T} \\
Q = Calor (J) \\
m = Massa (kg) \\
c = Calor específico (J/kg$\cdot$K) \\
$\Delta T$ = Variação de temperatura (K) \\
\vspace{1em}

% F = k q1 q2 / r^2
\textbf{F = k \frac{q_1 q_2}{r^2}} \\
F = Força elétrica (N) \\
k = 9,0 $\times$ 10$^9$ N m$^2$/C$^2$ \\
$q_1$, $q_2$ = Cargas (C) \\
r = Distância (m) \\
\vspace{1em}

\end{multicols}

\end{document}
